\begin{longtblr}
[
 label={tab:pico_air_quality},
 caption={Marco PICo utilizado para la revisión del estado del arte. Elaboración propia.}]{
colspec={X[j,m,0.7] X[j,m,3] X[j,m,5]},
width=\textwidth
}
\textbf{Compo- nente} & \textbf{Término resumido} & \textbf{Descripción operativa} \\

P &
Calidad del aire urbano &
Evaluada a partir de concentraciones de contaminantes atmosféricos (por ejemplo, PM$_{2.5}$, PM$_{10}$, NO$_2$, O$_3$) y variables ambientales asociadas, utilizadas como indicadores del estado de la calidad del aire. \\

I &
Lógica difusa para evaluación ambiental &
Aplicación de sistemas de inferencia basados en lógica difusa (por ejemplo, Mamdani o Sugeno) para la evaluación del nivel de calidad del aire y del riesgo asociado, mediante reglas lingüísticas y funciones de pertenencia. \\

Co &
Monitoreo en tiempo real &
Plataformas y sistemas de monitoreo en tiempo real orientados al análisis y visualización de datos de calidad del aire, provenientes de sensores ambientales y/o fuentes de datos abiertas. \\

\end{longtblr}
